\documentclass[runningheads]{llncs}
% \usepackage[pagebackref]{hyperref}

% ---------------------------------------------------------------
% Include basic ECCV package
 
% TODO REVIEW: Insert your submission number below by replacing '*****'
% TODO FINAL: Comment out the following line for the camera-ready version
% \usepackage[review,year=2026,ID=1259]{eccv}
% TODO FINAL: Un-comment the following line for the camera-ready version
\usepackage{eccv}

% OPTIONAL: Un-comment the following line for a version which is easier to read
% on small portrait-orientation screens (e.g., mobile phones, or beside other windows)
%\usepackage[mobile]{eccv}


% ---------------------------------------------------------------
% Other packages

% Commonly used abbreviations (\eg, \ie, \etc, \cf, \etal, etc.)
\usepackage{eccvabbrv}
% \usepackage[pagebackref]{hyperref}

% Include other packages here, before hyperref.
\usepackage{graphicx}
\usepackage{booktabs}

% The "axessiblity" package can be found at: https://ctan.org/pkg/axessibility?lang=en
\usepackage[accsupp]{axessibility}  % Improves PDF readability for those with disabilities.
\definecolor{myblue}{RGB}{193,248,255}

% ---------------------------------------------------------------
% Hyperref package

% It is strongly recommended to use hyperref, especially for the review version.
% Please disable hyperref *only* if you encounter grave issues.
% hyperref with option pagebackref eases the reviewers' job, but should be disabled for the final version.
%
% If you comment hyperref and then uncomment it, you should delete
% main.aux before re-running LaTeX.
% (Or just hit 'q' on the first LaTeX run, let it finish, and you
%  should be clear).

% TODO FINAL: Comment out the following line for the camera-ready version
\usepackage[pagebackref,breaklinks,colorlinks,citecolor=eccvblue]{hyperref}
% TODO FINAL: Un-comment the following line for the camera-ready version
% \usepackage{hyperref}
\usepackage{booktabs,multirow}
\usepackage{makecell}
\usepackage{array}
\usepackage[table]{xcolor}
\usepackage{xspace}
\usepackage{graphicx}
\usepackage{amsmath}
\usepackage{amssymb}
\usepackage{caption}
%\usepackage[compatibility=false]{caption}
\usepackage{subcaption}
% Support for ORCID icon
\usepackage{orcidlink}
\definecolor{colorNormal}{RGB}{193,248,255} % 浅蓝色 (Normal 最佳，与myblue一致)
\definecolor{colorDisco}{RGB}{255,235,205}  % 浅橙色 (Disco 最佳)
\definecolor{colorTwo}{RGB}{193,248,255} % 浅蓝色 (Single View 最佳)
\definecolor{colorSingle}{RGB}{255,235,205}  % 浅橙色 (Two Views 最佳)

\begin{document}

% ---------------------------------------------------------------
% TODO REVIEW: Replace with your title
\title{R3D: Revisiting 3D Policy Learning\vspace{-9pt}} 

% TODO REVIEW: If the paper title is too long for the running head, you can set
% an abbreviated paper title here. If not, comment out.
\titlerunning{R3D: Revisiting 3D Policy Learning}

% TODO FINAL: Replace with your author list. 
% Include the authors' OCRID for the camera-ready version, if at all possible.

\author{Zhengdong Hong*\inst{1}\thanks{Equal contribution. $\dagger$ Corresponding author.}\and
Shenrui Wu*\inst{2}\and
Haozhe Cui*\inst{2}\and
Boyi Zhao\inst{2}\and
Ran Ji\inst{2}\and
Yiyang He\inst{2}\and
Hangxing Zhang\inst{2}\and
Zundong Ke\inst{2}\and
Jun Wang\inst{1}\and
Guofeng Zhang\inst{1}\and
Jiayuan Gu$\dagger$\inst{2}
}
% \author{Zhengdong Hong\inst{1}\orcidlink{0000-1111-2222-3333} \and
% Second Author\inst{2,3}\orcidlink{1111-2222-3333-4444} \and
% Third Author\inst{3}\orcidlink{2222--3333-4444-5555}}

% TODO FINAL: Replace with an abbreviated list of authors.
\authorrunning{Z.Hong, S.Wu, H.Cui et al.}
% First names are abbreviated in the running head.
% If there are more than two authors, 'et al.' is used.
\institute{\vspace{-7pt}$^1$Zhejiang University \quad $^2$ShanghaiTech University}
% \institute{$^1$Zhejiang University \quad $^2$ShanghaiTech University \quad $^3$Univeristy of California, San Diego}
% TODO FINAL: Replace with your institution list.
% \institute{Princeton University, Princeton NJ 08544, USA \and
% Springer Heidelberg, Tiergartenstr.~17, 69121 Heidelberg, Germany
% \email{lncs@springer.com}\\
% \url{http://www.springer.com/gp/computer-science/lncs} \and
% ABC Institute, Rupert-Karls-University Heidelberg, Heidelberg, Germany\\
% \email{\{abc,lncs\}@uni-heidelberg.de}}

\maketitle

\vspace{-6mm}
\begin{abstract}
3D policy learning promises superior generalization and cross-embodiment transfer, but progress has been hindered by training instabilities and severe overfitting, precluding the adoption of powerful 3D perception models. In this work, we systematically diagnose these failures, identifying the omission of 3D data augmentation and the adverse effects of Batch Normalization as primary causes. We propose a new architecture coupling a scalable transformer-based 3D encoder with a diffusion decoder, engineered specifically for stability at scale and designed to leverage large-scale pre-training. Our approach significantly outperforms state-of-the-art 3D baselines on challenging manipulation benchmarks, establishing a new and robust foundation for scalable 3D imitation learning. Project Page: \url{https://r3d-policy.github.io/}
\vspace{-3mm}
  \keywords{3D Policy \and Imitation Learning}
\end{abstract}

\vspace{-6mm}
\section{Introduction}
\label{sec:intro}
\vspace{-1mm}
The landscape of robotic imitation learning has been significantly shaped by the success of 2D-centric policies, notably Diffusion Policy~\cite{chi2024diffusionpolicyvisuomotorpolicy} and ACT~\cite{zhao2023learning}, which introduced powerful paradigms like diffusion-based action decoding and action chunking. These techniques have been scaled up in modern Vision-Language-Action (VLA) models~\cite{zitkovich2023rt,kim2024openvla,black2024pi_0,bu2025univla}.
However, the reliance on 2D projections of an inherently 3D world introduces a fundamental limitation that hinders generalization to novel viewpoints~\cite{xie2024gengap} and complicates cross-embodiment skill transfer~\cite{o2024open}. In response, 3D policy learning~\cite{ze20243d,ke20243ddiffuseractorpolicy,ze2024idp3} is emerging as a compelling alternative. By reasoning over a canonical 3D scene representation, such policies can naturally fuse multi-view camera inputs, achieve inherent viewpoint invariance, and establish a consistent world frame that facilitates robust skill transfer across different robot morphologies.


Despite its promise, progress in 3D policy learning has been impeded by significant, often counterintuitive challenges. Pioneering studies like DP3~\cite{ze20243d} have identified a discrepancy in performance scaling, where a lightweight PointNet backbone outperformed more complex and powerful architectures. Furthermore, as shown in section~\ref{fix_overfitting_issue} and section~\ref{sec:ablation}, training 3D policies can be quite unstable, especially when model capacity increases. These two issues—\textbf{unexpected scaling and unstable training}—severely hinder the adoption of modern 3D architectures~\cite{wu2024pointtransformerv3,zhou2023uni3d,zhou2024point} for better 3D policy learning.

In this paper, we systematically address these challenges to unlock the potential of 3D policy learning. We begin with a comprehensive analysis of existing frameworks, identifying and remedying instabilities that have previously limited performance. We reveal two crucial oversights. First, the common implementation of DP3~\cite{ze20243d} omits \textbf{data augmentation}, a standard technique for stabilizing training and mitigating overfitting. Second, we identify that \textbf{Batch Normalization} significantly degrades the performance of more powerful 3D backbones.

Motivated by these insights, we propose a novel architecture and a set of recipes for scaling 3D policy learning. Our approach integrates a high-capacity, \textbf{LayerNorm-only 3D encoder} inspired by recent advances in 3D segmentation~\cite{zhou2024point}, coupled with a Diffusion Transformer~\cite{peebles2023scalable} for action decoding. Crucially, our architecture \textbf{preserves the spatial resolution of 3D features} rather than collapsing them into a global vector. These designs allow the model to fully leverage scaled-up 3D encoders while ensuring training stability, consistent inference behavior, and superior spatial reasoning. Besides, such designs also enable \textbf{pretraining} of the 3D encoder on 3D segmentation tasks, which accelerates convergence and boosts final performance. Finally, we introduce an \textbf{auxiliary task} to jointly decode end-effector poses and joint angles; by maintaining spatial resolution throughout the decoding process, this task enhances the model’s proprioceptive awareness and joint control precision.


Our main contributions are as follows:
\begin{itemize}
    \item We identify critical training practices—\textbf{3D data augmentation} and \textbf{removing Batch Normalization}—that enable the stable training of large-scale 3D policies and prevent the overfitting issues.
    \item We propose a \textbf{novel 3D policy architecture} that integrates a scalable transformer-based point cloud encoder with a diffusion-based action decoder. This design explicitly enables \textbf{pre-training} of the encoder on large-scale 3D segmentation tasks to facilitate knowledge transfer.
    \item We conduct extensive experiments showing that our model \textbf{significantly outperforms} state-of-the-art 3D imitation learning baselines on both simulated and real-world manipulation tasks. We also analyze the impact of \textbf{encoder capacity}, \textbf{decoder depth}, and \textbf{camera configurations}, providing valuable empirical references for designing future 3D policies.
\end{itemize}


\section{Related Work}

\label{sec:related}

% \subsection{2D Perception and Representation Learning}
\subsection{2D Policy Learning}
% \vspace{-0.2em}
Early visuomotor policy learning primarily relied on 2D visual inputs, focusing on imitation and end-to-end action prediction from RGB observations. Diffusion-based approaches such as Diffusion Policy~\cite{chi2024diffusionpolicyvisuomotorpolicy},  RDT-1B~\cite{liu2025rdt1bdiffusionfoundationmodel}, and CF-SDP~\cite{yu2025efficienttaskspecificconditionaldiffusion} demonstrate stable behavior cloning through generative modeling in pixel space, while transformer-based architectures like ACT~\cite{zhao2023learning} and Perceiver-Actor~\cite{shridhar2022perceiveractormultitasktransformerrobotic} leverage spatiotemporal attention for long-horizon manipulation.  
In parallel, recent 2D VLA models further extend visuomotor learning to multimodal settings. Systems such as OpenVLA~\cite{kim2024openvla}, OpenVLA-OFT~\cite{kim2025fine}, $\pi_0$~\cite{black2024pi_0}, RT-2~\cite{zitkovich2023rt}, RT-X~\cite{o2024open}, RoboFlamingo~\cite{li2023vision}, Octo~\cite{team2024octo}, GR-1~\cite{wu2023unleashing}, and UniVLA~\cite{bu2025univla} integrate large vision–language backbones with robot action policies, enabling semantic grounding and instruction following via language-conditioned control. 
Despite the rapid progress of these 2D models, their reliance on image-space representations inherently limits spatial understanding. Without explicit depth or 3D cues, they often exhibit weak geometric reasoning, poor viewpoint generalization, and reduced robustness in cluttered or spatially complex environments. These limitations motivate the development of depth-aware and geometry-aware formulations that bridge 2D perception with 3D understanding.

% \subsection{Bridging 2D Perception and 3D Understanding}
\vspace{-0.06em}
\subsection{Policy Learning with Implicit 3D}
To overcome the geometric limitations of 2D inputs, recent works have explored a spectrum of depth- and multi-view–aware representations that bridge 2D perception with 3D understanding. A first line of approaches focuses on deriving implicit or lifted geometric structure from monocular RGB observations. GP3~\cite{qian2025gp33dgeometryawarepolicy} constructs geometry-aware policies via multi-view fusion, while VO-DP~\cite{ni2025vodpsemanticgeometricadaptivediffusion} leverages semantic–geometric cues to achieve implicit depth reasoning without explicit 3D supervision. Lift3D~\cite{jia2024lift3dfoundationpolicylifting} further lifts large-scale 2D pretrained models into a 3D-aligned feature space, enhancing robustness under spatial variations. Leveraging multi-view consistency, 3D-MVP~\cite{qian20253dmvp3dmultiviewpretraining} and FVP~\cite{Hou_2025_ICCV} introduce multi-view pretraining strategies that enable visual backbones to learn geometry-aware representations from synthetic or real multi-view data.  
Meanwhile, VLA frameworks such as DepthVLA~\cite{yuan2025depthvlaenhancingvisionlanguageactionmodels}, SpatialVLA~\cite{qu2025spatialvlaexploringspatialrepresentations}, 3D-CAVLA~\cite{bhat20253d}, and Spatial Forcing~\cite{li2025spatialforcingimplicitspatial} introduce depth cues or egocentric spatial encodings to improve geometric grounding within multimodal policy models. These methods constitute the emerging 2.5D regime—mitigating projection artifacts and enhancing geometric reasoning, yet still limited by partial 3D awareness and incomplete spatial consistency between perception and action.

% \subsection{3D Policy and Geometric Reasoning}
\vspace{-0.06em}
\subsection{3D Policy Learning}
Recent advances in 3D policy learning encompass diffusion-based, transformer-based, flow-based, and VLA formulations. 
Diffusion-based approaches such as 3D Diffusion Policy (DP3)~\cite{ze20243d}, 3D Diffuser Actor~\cite{ke20243ddiffuseractorpolicy}, AC-DiT~\cite{chen2025acditadaptivecoordinationdiffusion}, and AnchorDP3~\cite{zhao2025anchordp33daffordanceguided} incorporate point-cloud or volumetric encodings into transformer-augmented diffusion frameworks to enhance spatial reasoning and generalization. 
Transformer-based architectures, including FP3~\cite{yang2025fp33dfoundationpolicy}, Act3D~\cite{gervet2023act3d3dfeaturefield}, and RVT~\cite{pmlr-v229-goyal23a}, emphasize scalable multi-view representation learning and 3D spatial aggregation for manipulation. 
Recently, studies extend 3D policy learning into the VLA domain, such as 3D-VLA~\cite{zhen20243dvla3dvisionlanguageactiongenerative}, GeoVLA~\cite{sun2025geovlaempowering3drepresentations}, PointVLA~\cite{li2025pointvlainjecting3dworld}, and BridgeVLA~\cite{li2025bridgevlainputoutputalignmentefficient}.
They integrate 3D geometric or depth-aware spatial representations into large VLA frameworks.
Flow-based formulations such as 3D FlowMatch Actor~\cite{gkanatsios20253dflowmatchactorunified} and ManiFlow~\cite{yan2025maniflowgeneralrobotmanipulation} combine diffusion and flow mechanisms to improve motion smoothness and sample efficiency.  
Despite these advances, current 3D policies still face challenges in training stability, computational cost of dense 3D encoding, and redundant 3D representations—motivating our work’s focus on a lean transformer-based 3D diffusion formulation with scalable training.


\section{Revisiting 3D Policy Learning}
\label{sec:revisiting}

To enable the scaling of 3D policy learning, we first conduct a systematic investigation into why prior methods fail to benefit from more powerful 3D backbones. Our analysis reveals a critical insight: performance is dictated not by architectural novelty, but by several subtle yet decisive implementation details. In this section, we present these findings and distill them into a set of actionable training principles designed to guide future research.

\noindent\textbf{Experimental Setup.}
We utilize RoboTwin 2.0~\cite{chen2025robotwin}, a simulation benchmark, as our testbed to revisit these design choices. We select five representative bimanual tasks and report success rates as our primary metric. Additionally, we provide learning curves to evaluate learning stability. For comprehensive experimental settings, please refer to section~\ref{sec: Exp}.

\subsection{Overcoming the Scaling Paradox: From BN to LN}
Prior work such as DP3~\cite{ze20243d} suggests that lightweight PointNet~\cite{qi2017pointnet} backbones perform better than more powerful ones. To investigate this ``scaling paradox'', we replace the PointNet backbone with a high-capacity encoder used in Uni3D~\cite{zhou2023uni3d}. This architecture employs Farthest Point Sampling (FPS) to sample center points and $k$-Nearest Neighbors ($k$-NN) to extract neighborhood patches, which are then processed by a Transformer initialized from 2D Vision Transformers~\cite{dosovitskiy2020image}.

Our initial experiments confirm the findings in DP3: naively adopting this stronger backbone leads to a significant performance degradation. However, further investigation identifies \emph{Batch Normalization} (BN) as the primary culprit. BN layers often struggle with the high variance and small batch sizes typical of imitation learning. 

To validate this insight, we compare four configurations: DP3 with PointNet (BN vs. LN) and DP3 with Uni3D (BN vs. LN). As shown in Tab.~\ref{tab:newtab0}, replacing BN with \emph{Layer Normalization} (LN) allows the Uni3D backbone to not only become trainable but to significantly outperform PointNet. Notably, LN doesn't degrade PointNet's performance, suggesting it's a more robust default for 3D policies.
% \todo{Supported with Table. Compare: 1. DP3(PointNet,BN), 2. DP3(PointNet,LN), 3. DP3(Uni3D,LN), 4. DP3(Uni3D,BN) 5. fp3(BN)}.
%fp3仍然没有注意到 DP3 BN not work 1,3必须出现 核心结论：去掉BN，更大网络可以work
% \vspace{-1.5em}
\begin{table*}[]
\centering
\small
\captionsetup{justification=centering}

% --- 左侧：Table X (BN vs. LN) ---
\begin{minipage}[]{0.48\textwidth} % 稍微缩小宽度比例确保不重叠
% \vspace{-1em}
    \centering
    \caption{\textbf{Lessons in BN vs. LN.} \\We use DP3\cite{ze20243d} default hyperparameters for training.}
    % \vspace{-1em}
    \label{tab:newtab0}
    % \vspace{2pt} % 统一标题与表格的间距
    \resizebox{\linewidth}{!}{
        % 通过大幅增加 arraystretch 来补偿行数差（4行 vs 6行）
        % 这里的 1.55 是根据 (6行/4行) 的比例并考虑表头高度微调得出的
        \renewcommand{\arraystretch}{1.55} 
        \begin{tabular}{l|ccc|c}
            \toprule
            Method & \shortstack{Beat \\ Hammer} & \shortstack{Move \\ Card} & \shortstack{Place \\ Shoe} & \textbf{Avg.} \\
            \midrule
            DP3 (PointNet+BN) & 0 & 3 & 0 & 1.0 \\
            DP3 (PointNet+LN) & 79 & 57 & 43 & 59.6 \\
            DP3 (Uni3D+BN)     & 0 & 0 & 0  & 0.0 \\
            DP3 (Uni3D+LN)     & 86 & 60 & 48 & 64.7 \\
            \bottomrule
        \end{tabular}
    }
\end{minipage}
\hfill % 必须保留，用于撑开左右间距
% --- 右侧：Table Y (Encoder Sizes) ---
\begin{minipage}[]{0.48\textwidth}
    \centering
    \caption{Different encoder sizes (using RoboTwin 2.0~\cite{chen2025robotwin} default 1024 points).}
    % \vspace{-1em}
    \label{tab:encoder_size_comparison}
    % \vspace{2pt}
    \resizebox{\linewidth}{!}{
        \renewcommand{\arraystretch}{1.1} % 右表保持紧凑
        \begin{tabular}{l ccc}
            \toprule
            Tasks & \shortstack[c]{tiny \\ (53MB)} & \shortstack[c]{small \\ (115MB)} & \shortstack[c]{base \\ (366MB)} \\
            \midrule
            Move Playingcard   & 85 & \cellcolor{myblue}\textbf{92} & 83 \\
            Turn Switch        & \cellcolor{myblue}\textbf{56} & 55 & 50 \\
            Place Shoe         & \cellcolor{myblue}\textbf{87} & 68 & 70 \\
            Open Microwave     & \cellcolor{myblue}\textbf{100} & 93 & 77 \\
            Beat Block Hammer  & \cellcolor{myblue}\textbf{91} & 76 & 85 \\
            \midrule
            \textbf{Average}   & \cellcolor{myblue}\textbf{83.8} & 76.8 & 73.0 \\
            \bottomrule
        \end{tabular}
    }
\end{minipage}
% \vspace{-10pt}
\end{table*}
% \begin{table}[]
%     \centering
%     \small
%     \renewcommand{\arraystretch}{1.15}
%     \setlength{\tabcolsep}{3pt}

%     \caption{\textbf{Lessons in BN vs. LN.} Models marked with $^*$ are trained using DP3\cite{ze20243d} hyperparameters.}
%     \label{tab:newtab0}
%     \resizebox{\linewidth}{!}{
%     \begin{tabular}{l|ccc|c}
%         \toprule
%         Method &
%         \shortstack{Beat Block\\
% Hammer} &
%         \shortstack{Move Playingcard\\
% Away} &
%         \shortstack{Place\\
% shoe} &
%         \textbf{Average} \\
%         \midrule
%         % DP3 (PointNet+BN+nocolor)           & 0 & 0 &  0 & 0 \\
%         % DP3 (PointNet+LN+nocolor)              & 89 & 71 &  53 & 71 \\
%         % DP3 (PointNet+BN)          & 0 & 0 &  0 & 0 \\
%         % DP3 (PointNet+LN)              & 82 & 59 &  37 & 59.3 \\
%         *DP3 (PointNet+BN)           & 0 & 3 &  0 & 1 \\    %dp3 parameters
%         *DP3 (PointNet+LN)              & 79 & 57 &  43 & 59.6 \\     %%dp3 parameters
%         DP3 (Uni3D+BN)              & 0 & 0 & 0  & 0 \\
%         DP3 (Uni3D+LN)              & 86 & 60 & 48  & 64.7 \\
%         % *DP3 (Uni3D+LN+nocolor)              & 72 & 78 & 68  & 72.7 \\
%         % FP3 (Uni3D+BN) &  &  &   &  \\
%         % ACT              &  &  &   &  \\
%         % \midrule
%         % \textbf{Ours}    & \textbf{82} & \textbf{64} & \textbf{12} & \textbf{52.67} \\
%         \bottomrule
%     \end{tabular}}
% \end{table}
\vspace{-2em}
\subsection{Mitigating Overfitting through Data Augmentation}\label{fix_overfitting_issue}
3D policy training is notoriously vulnerable to overfitting~\cite{ze20243d}, often compelling researchers to report peak performance over a range of checkpoints rather than final convergence. We address it by introducing a data augmentation pipeline consisting of three key strategies.
% \begin{figure}
%     \centering
%     \includegraphics[width=0.4\textwidth]{figures/data-aug131_cropped.pdf}
%     \caption{Success Rate Differences Using Data Augmentation}
%     \label{fig:data-aug}
% \end{figure}
\textbf{1. FPS Randomization:}
Following PointNet++~\cite{qi2017pointnet++}, we randomize the point order in each training step for backbones utilizing FPS. This prevents the model from relying on the deterministic order induced by the sampling algorithm, making the encoder learn more robust spatial representations. \textbf{2. Point Cloud Color Jitter:} Inspired by image augmentations in Qt-Opt~\cite{kalashnikov2018qtoptscalabledeepreinforcement}, we apply color jittering directly to the point cloud's RGB channels. We introduce random variations in brightness (range [$-0.125, 0.125$]), contrast (range [$0.5, 1.5$]), and saturation (range [$0.5, 1.5$]). \textbf{3. Additive Noise and Dropout:} We apply modest Gaussian noise to both point cloud coordinates and proprioceptive vectors. Furthermore, we employ random point dropout, forcing the network to maintain accurate predictions even from incomplete scene representations.

The impact of these augmentations is evidenced by our training and validation curves. In Fig.~\ref{fig:data-aug}, we observe that without augmentation, there are significant fluctuations in the learning curves and a marked decline in success rates as training progresses. The introduction of data augmentation not only stabilized the training process but also significantly improved success rates. This improvement arises because the policy leverages a broader range of out-of-distribution data, thereby enhancing its generalization capabilities.
% \todo{Supported with training/validation curves}
% 补充实验，base:Uni3d/Poin    tNet 

% \subsection{Enhancing Training Efficiency by Action Filtering} 
% We observe that demonstration datasets often contain 'static actions' where joint velocities are near zero. Since removing these frames does not affect training outcomes but improves computational efficiency, we applied action filtering to the data.
% \todo{maybe add some results to support this add action filtering(Based on B)}
% \begin{figure}
%     \centering
%     \includegraphics[width=0.7\linewidth]{example-image}
%     \caption{data efficiency by using action filtering}
%     \label{fig:action-filter}
% \end{figure}

% \begin{figure}
%     \centering
%     \includegraphics[width=0.85\linewidth]{figures/decoder_depth_ablation_final_cropped.pdf}
%     \caption{Success Rates with Varied Decoder Depth.}
%     \label{fig:decoder_depth}
% \end{figure}
\begin{figure*}[t] % 使用 figure* 确保在双栏模式下也能正常显示
    \centering
    
    % --- 第一张图 (Figure X) ---
    \begin{minipage}{0.47\textwidth}
        \centering
        \includegraphics[width=0.9\textwidth]{figures/data-aug131_cropped.pdf}
        \caption{Comparison of learning curves with vs. without data augmentation.}
        \label{fig:data-aug}
    \end{minipage}
    \hfill % 撑开中间的间距
    % --- 第二张图 (Figure X+1) ---
    \begin{minipage}{0.50\textwidth}
        \centering
        \includegraphics[width=1\textwidth]{figures/decoder_depth_ablation_final-2_cropped.pdf}
        \caption{Success rates with varied decoder depth (each point is averaged over 5 ckpts).}
        \label{fig:decoder_depth}
    \end{minipage}
\vspace{-10pt}
\end{figure*}
%，效率上可能有提升
% \begin{table}[]
%     \centering
%     \small
%     \renewcommand{\arraystretch}{1.15}
%     \setlength{\tabcolsep}{3pt}

%     \caption{\textbf{Lessons learning from using data-augmentation} All methods use pointcloud colors. 750 epochs}
%     \label{tab:newtab1}
%     \resizebox{\linewidth}{!}{
%     \begin{tabular}{l|ccc|c}
%         \toprule
%         Method &
%         \shortstack{Beat Block\\
% Hammer} &
%         \shortstack{Move Playingcard\\
% Away} &
%         \shortstack{Place\\
% shoe} &
%         \textbf{Average} \\
%         \midrule
     
%         DP3 (Uni3D+LN)              & 0.86 & 0.6 &  0.48 & 0.65 \\
%         DP3 (Uni3D+LN+data aug.)           & 0.90 & 0.60 &  0.60 & 0.70 \\
%         % DP3 (Uni3D+LN+action filt.)              &  &  &   &  \\
%         % DP3 (Uni3D+LN+data aug.+action filt.)              & 0.94 & 0.62 & 0.52  & 0.69 \\
%         % FP3 (Uni3D+BN) &  &  &   &  \\
%         % ACT              &  &  &   &  \\
%         % \midrule
%         % \textbf{Ours}    & \textbf{82} & \textbf{64} & \textbf{12} & \textbf{52.67} \\
%         \bottomrule
%     \end{tabular}}
% \end{table}
\section{Improving 3D Policy Architecture}
\label{sec:architecture}
\subsection{Architecture Inspired by Advanced 3D Learning}
% \begin{table}[]
%     \centering
%     \small
%     \renewcommand{\arraystretch}{1.15}
%     \setlength{\tabcolsep}{3pt}

%     \caption{\textbf{scaling the models}}
%     \label{tab:scaling-model}
%     \resizebox{\linewidth}{!}{
%     \begin{tabular}{l|ccc|c}
%         \toprule
%         Method &
%         \shortstack{Beat Block\\
% Hammer} &
%         \shortstack{Move Playingcard\\
% Away} &
%         \shortstack{Place\\
% shoe} &
%         \textbf{Average} \\
%         \midrule
     
%         Ours(pointsam-tiny,depth4, 7200)              &  &  &   &  \\
%         Ours(pointsam-small, 7210)           &  &  &   &  \\
%         Ours(pointsam-base, 7211)              &  &  &   &  \\
%         Ours(pointsam-tiny, depth2, 7220)              &  &  &   &  \\
%         Ours(pointsam-tiny, depth6, 7221)              &  &  &   &  \\
%         Ours(pointsam-tiny, depth8, 7222)              &  &  &   &  \\
%         Ours(ptv3-tiny, depth4, 7230)              &  &  &   &  \\
%         Ours(ptv3-small, depth4, 7231)              &  &  &   &  \\
%         % FP3 (Uni3D+BN) &  &  &   &  \\
%         % ACT              &  &  &   &  \\
%         % \midrule
%         % \textbf{Ours}    & \textbf{82} & \textbf{64} & \textbf{12} & \textbf{52.67} \\
%         \bottomrule
%     \end{tabular}}
% \end{table}

%scaling the data
% \begin{table}[]
%     \centering
%     \small
%     \renewcommand{\arraystretch}{1.15}
%     \setlength{\tabcolsep}{3pt}

%     \caption{\textbf{scaling the data}}
%     \label{tab:scale-data}
%     \resizebox{\linewidth}{!}{
%     \begin{tabular}{l|ccc}
%         \toprule
%         Method &
%         \shortstack{single} &
%         \shortstack{5-task} &
%         \shortstack{10-task}  \\
%         \midrule
     
%         \shortstack{Beat Block Hammer}              &  &  &   \\
%         \shortstack{Move Playingcard Away}           &  &  &     \\
%         \shortstack{Turn Switch}              &  &  &     \\
%         \shortstack{Open Microwave}              &  &  &   \\
%         \shortstack{Place Shoe}              &  &  &     \\
%         % FP3 (Uni3D+BN) &  &  &   &  \\
%         % ACT              &  &  &   &  \\
%         % \midrule
%         % \textbf{Ours}    & \textbf{82} & \textbf{64} & \textbf{12} & \textbf{52.67} \\
%         \bottomrule
%     \end{tabular}}
% \end{table}
%\shortstack{Beat Block \\ Hammer} &
%\shortstack{Move Playingcard \\ Away} &
%\shortstack{Turn \\ Switch} &
%\shortstack{Open \\ Microwave} &
%\shortstack{Place \\ shoe} &

% \begin{figure*}[htbp]
% \centering
% \includegraphics[width=\textwidth]{figures/old-pipeline_cropped.pdf}% Reduce the figure size so that it is slightly narrower than the column.
% \caption{Pipeline Architecture: our 3D policy architecture consists of a point cloud encoder encoding input point clouds and a diffusion transformer to decode actions. The encoder can be pre-trained on 3D segmentation tasks.}
% %ee和joint用不同颜色来画token的样子
% \label{fig:pipe}
% \end{figure*}

\begin{figure*}[t]
\centering
\includegraphics[width=\textwidth]{figures/R3D_pipeline-3_cropped.pdf}% Reduce the figure size so that it is slightly narrower than the column.
\caption{\textbf{Pipeline architecture}: our 3D policy architecture consists of a point cloud encoder encoding input point clouds and a diffusion transformer to decode actions. The encoder can be pre-trained on 3D segmentation tasks.}
%ee和joint用不同颜色来画token的样子
\label{fig:pipe-overall}
\end{figure*}
% \vspace{-4.5em}

\begin{figure*}[t]
\centering
\includegraphics[width=\textwidth]{figures/figure-1-1-131-new_cropped.pdf}% Reduce the figure size so that it is slightly narrower than the column.
\caption{Difference between decoders used in DP3~\cite{ze20243d} and ours. DP3 uses a global feature extracted from the input point cloud, and conditions the denoising process via FiLM. Ours uses dense features, and conditions the denoising process through cross-attention. Figure b) details the structure of a single attention block.}
\label{fig:pipe}
\vspace{-15pt}
\end{figure*}

We propose a unified perspective on 3D perception and 3D policy learning, observing that both domains rely on the effective interaction between geometric features and task-specific queries. In 3D segmentation, for instance, architectures like Point-SAM~\cite{zhou2024point} utilize a two-way Transformer to iteratively update point cloud tokens and prompt tokens via cross-attention. This process enriches prompt tokens with geometric context while conditioning point features on prompt-specific information to highlight regions of interest. Diffusion-based policies~\cite{chi2024diffusionpolicyvisuomotorpolicy} share a  similar underlying logic: observation tokens are cross-attended to noisy action queries, which are then iteratively denoised. 

Leveraging this architectural synergy, we design a scalable 3D policy that adapts proven perception backbones for robotic control. In the following sections, we detail our high-capacity point cloud encoder (section~\ref{sec:encoder}), our spatially-aware Diffusion Transformer action decoder (section~\ref{sec:decoder}), and an auxiliary task (section~\ref{sec:auxiliary}) designed to enhance proprioceptive grounding. Fig.~\ref{fig:pipe-overall} illustrates the complete architecture.

% \usepackage{tabularx} % 必须在导言区加上此包


% \begin{table}[htbp]
%     \centering
%     \small 
%     % \captionsetup{justification=centering} % 如果全局已配置可省略
%     \caption{Success Rates Across Different Encoder Sizes}
%     \label{tab:encoder_size_comparison}
%     \begin{tabular*}{\linewidth}{l @{\extracolsep{\fill}} c c c} 
        
%         \toprule
%         Tasks & 
%         \shortstack[c]{tiny \\ (53MB)} & 
%         \shortstack[c]{small \\ (115MB)} & 
%         \shortstack[c]{base \\ (366MB)} \\
        
%         \midrule
        
%         \makecell[l]{Move Playingcard} & \makecell{85} & \makecell{\textbf{92}} & \makecell{83} \\
        
%         \makecell[l]{Turn Switch} & \makecell{\textbf{56}} & \makecell{55} & \makecell{50} \\
        
%         \makecell[l]{Place Shoe} & \makecell{\textbf{87}} & \makecell{68} & \makecell{70} \\
        
%         \makecell[l]{Open Microwave} & \makecell{\textbf{100}} & \makecell{93} & \makecell{77} \\
        
%         \makecell[l]{Beat Block Hammer} & \makecell{\textbf{91}} & \makecell{76} & \makecell{85} \\
        
%         \midrule
        
%         \makecell[l]{Average} & \makecell{\textbf{83.8}} & \makecell{76.8} & \makecell{73.0} \\
        
%         \bottomrule
        
%     \end{tabular*} 
%     \vspace{-2em}
% \end{table}
% \noindent\textbf{Notation.}
% We define $T_o$ as the observation horizon and $T_a$ as the action prediction horizon. At each timestep $t$, the observation consists of a point cloud $P_t \in \mathbb{R}^{N_P \times 3}$ and a proprioceptive state $q_t \in \mathbb{R}^{N_q}$, where $N_P$ is the number of points and $N_q$ is the dimensionality of the robot's configuration space. Input point clouds are sampled or padded to a fixed size $N_P$.
% \todo{Change into a preliminary subsection, also explaining diffusion?}

\subsection{Preliminaries}
\label{sec:preliminaries}
\noindent\textbf{Problem Setup.} We define $T_o$ as the observation horizon and $T_a$ as the action prediction horizon. At each timestep $t$, the observation $O_t$ consists of a point cloud $P_t \in \mathbb{R}^{N_P \times 6}$ (with point colors) and a proprioceptive state $q_t \in \mathbb{R}^{N_q}$, where $N_P$ is the number of points and $N_q$ is the dimensionality of the robot's configuration space. Input point clouds are sampled or padded to a fixed size $N_P$. The goal is to learn a policy that predicts a sequence of actions $\mathbf{a}_t = \{a_t, a_{t+1}, \dots, a_{t+T_a-1}\}$.

\noindent\textbf{Diffusion Formulation.} We model the policy using a conditional diffusion process. Following~\cite{chi2024diffusionpolicyvisuomotorpolicy}, the model learns to refine a noise sequence $\mathbf{a}^{(K)} \sim \mathcal{N}(0, \mathbf{I})$ into a clean action sequence $\mathbf{a}^{(0)}$ over $K$ iterations. At each iteration $k$, the model predicts the noise or the denoised action conditioned on the observation $O_{t-T_o+1:t}$. In our work, we employ a Diffusion Transformer (DiT) as the noise prediction network $\epsilon_\theta$.

\subsection{High-Capacity 3D Geometric Encoder}
\label{sec:encoder}
Our architecture utilizes a powerful 3D encoder inspired by recent advances in point cloud segmentation~\cite{zhou2024point}. Following the insights from section~\ref{sec:revisiting}, the encoder is a Transformer-based backbone that exclusively uses Layer Normalization (LN) to ensure scaling stability. 

The encoding process begins by partitioning the input point cloud $\mathcal{P} \in \mathbb{R}^{N_p \times 3}$ into a set of local patches. Specifically, we use Farthest Point Sampling (FPS) to select $N_C$ center points and $k$-Nearest Neighbors ($k$-NN) to group their local neighborhoods. Each patch is then projected into a $D$-dimensional latent embedding via a lightweight PointNet~\cite{qi2017pointnet} with LN. These patch embeddings are augmented with positional embeddings of patch centers and then processed through a series of self-attention layers initialized from a pre-trained 2D Vision Transformer (ViT)~\cite{dosovitskiy2020image}.

Crucially, our encoder outputs a set of $N_C$ structured point tokens, $F^{geo} \in \mathbb{R}^{N_C \times D}$, representing the fine-grained geometry of the scene. This stands in contrast to prior work DP3~\cite{ze20243d} where features are typically collapsed into a single global descriptor. Even in similar architectures like Uni3D~\cite{zhou2023uni3d}, a $[CLS]$ token is usually appended to extract a global feature. We intentionally omit this global pooling, instead preserving the full spatial resolution of the patch tokens to provide the downstream decoder with a dense geometric representation.

\subsection{Pretraining Encoder on 3D Segmentation Tasks}
Besides, since we borrow the encoder structure from 3D segmentation~\cite{zhou2024point}, we are able to pretrain our encoder on large-scale 3D scene datasets like ScanNet\cite{dai2017scannet}, ARKitScenes\cite{baruch2021arkitscenes}, and PartNeXt\cite{wangpartnext}. Pre-training the 3D visual encoder serves two primary purposes. First, it injects rich, generalizable geometric priors that substantially reduce the number of task-specific samples required when fine-tuning on robot-centric datasets. Second, pre-training on diverse sources improves robustness to distributional shifts encountered in real-world robotic environments (e.g., sensor noise, partial observations, varying scene complexity). 

% \todo{adding pretraining here}
% \todo{key components are here by Prof. Gu, other details should be checked and added. Some might be missing.}
% 我们这里可以简单提一句，这样的设计可以enable我们可以做pretraining on large 3d dataset（e.g.），后面在实验里面可以在setting里面讲我们在什么dataset上pretraining，以及一些必要的细节，以及实验结果分析中可以再次突出pretraining

\subsection{Spatially-Aware Action Decoding}
\label{sec:decoder}

The action decoder is formulated as a Diffusion Transformer~\cite{peebles2023scalable} that operates on a sequence of noisy actions $\mathbf{a}^{(k)}_t = (a^{(k)}_t, a^{(k)}_{t+1}, \dots, a^{(k)}_{t+T_a-1})$ at diffusion iteration $k$. Rather than conditioning the denoising process on a global latent vector, our decoder performs cross-attention directly between the action queries and the full set of geometric tokens.

The decoder processes three types of tokens, all projected into a shared $D$-dimensional space:
\begin{enumerate}
    \item \textbf{Geometric Tokens:} The historical sequence of point cloud tokens from the observation horizon, $\mathcal{F}^{geo} = \{F^{geo}_{t-T_o+1}, \dots, F^{geo}_t\} \in \mathbb{R}^{(T_o \cdot N_P) \times D}$.
    \item \textbf{Proprioception Tokens:} The robot states $\mathcal{F}^{prio} = \{F^{prio}_{t-T_o+1}, \dots, F^{prio}_t\} \in \mathbb{R}^{T_o \times D}$, where each $F^{prio}_t$ is generated by an MLP from $q_t$.
    \item \textbf{Action Tokens:} The noisy action sequence $\mathcal{F}^{act} = \{F^{act}_{t}, \dots, F^{act}_{t+T_a-1}\} \in \mathbb{R}^{T_a \times D}$, where each token $F^{act}_t$ is embedded from $a^{(k)}_t$ via an MLP.
\end{enumerate}

To incorporate temporal context, we add learnable temporal embeddings to all tokens based on their relative timestep within $T_o$ or $T_a$. Additionally, a sinusoidal embedding representing the diffusion iteration $k$ is prepended to the action sequence. The DiT architecture uses the observation tokens ($\mathcal{F}^{geo}, \mathcal{F}^{prio}$) as the context for cross-attention, while the action tokens $\mathcal{F}^{act}$ serve as the primary queries to be denoised. The final updated action tokens are decoded into action chunks via an MLP.

A key contribution of this design is the \textbf{preservation of spatial resolution}. By maintaining the spatial structure of the 3D features throughout the decoding process, the model retains the precise geometric relationships between the robot and the environment. This avoids the information bottlenecks inherent in global pooling, allowing the Diffusion Transformer to attend to specific task-relevant regions—such as a handle or a container edge—during the denoising iterations. This results in significantly more precise action predictions, particularly for tasks requiring fine-grained spatial reasoning. 
% \todo{Add a visualization of attention map.}

\subsection{Multi-Objective Decoding for Proprioceptive Grounding}
\label{sec:auxiliary}
To further regularize the learning process and improve control accuracy, we introduce an auxiliary task to decode task-space actions (end-effector poses) along with joint-space actions. Since the end-effector pose is inherently linked to the joint configuration via forward kinematics, this auxiliary objective acts as a ``free lunch'' that provides additional supervision without extra data collection.

Specifically, we include an auxiliary set of action tokens $^{ee}\mathcal{F}^{act} = \{^{ee}F^{act}_{t}, \dots,\\ ^{ee}F^{act}_{t+T_a-1}\} \in \mathbb{R}^{T_a \times D}$, embedded from noisy 7D task-space actions, consisting of 3D translation ($xyz$) and 4D quaternions ($wxyz$) represented in the base frame. During the diffusion process, the model simultaneously denoises both joint-space and task-space tokens. Notably, we employ a \emph{causal attention mask} within the self-attention layers of the decoder: while task-space tokens \emph{attend to} joint-space tokens, the joint-space tokens do not attend to the task-space ones. This unidirectional information flow ensures that the auxiliary task provides features for proprioceptive grounding without interfering with or degrading the primary joint-control objective. 

By leveraging the preserved spatial features, the model learns to ground its proprioception within the 3D scene. This multi-objective optimization encourages the network to develop a more robust internal representation of the robot's state relative to the environment, ultimately leading to higher joint control precision and better task-space alignment.
\subsection{Network Scaling in Policy Learning}
\label{sec: scaling}
\begin{table}[b]
    \centering
    \small
    \renewcommand{\arraystretch}{1.2} % 稍微增加行高，方便左侧内容居中对齐时更美观
    \setlength{\tabcolsep}{4pt}

    % 左侧表格：由于列数少，强制填满 0.48\textwidth 会显得更大
    \begin{minipage}[t]{0.48\textwidth}
        \centering
        \caption{Encoder size and number of demonstrations used in training (averaged on 5 RoboTwin 2.0~\cite{chen2025robotwin} ``Easy'' tasks).}
        \label{tab:supptab0}
        \resizebox{\linewidth}{!}{
            \begin{tabular}{l|ccc}
                \toprule
                Encoder Size & 50 demos & 100 demos & 200 demos \\
                \midrule
                \makecell[l]{~~~~~~~~Ours \\ (PointSAM-tiny)}  & \cellcolor{myblue}\textbf{83.8} & \cellcolor{myblue}\textbf{95.3} & \cellcolor{myblue}\textbf{97.8} \\
                \makecell[l]{~~~~~~~~Ours \\ (PointSAM-small)} & 76.8 & 89.0 & 92.4 \\
                \makecell[l]{~~~~~~~~Ours \\ (PointSAM-base)}  & 73.0 & 86.5 & 88.2 \\
                \bottomrule
                % 补偿右表多出的行数高度
            \end{tabular}
        }
    \end{minipage}
    \hfill
    % 右侧表格：严格等宽
    \begin{minipage}[t]{0.48\textwidth}
        \centering
        \caption{Encoder size (8192 points).}
        \label{tab:supptab1}
        \resizebox{\linewidth}{!}{
            \begin{tabular}{l|ccc}
                \toprule
                \shortstack[l]{Encoder Size} & \shortstack[c]{tiny \\ (53MB)} & \shortstack[c]{small \\ (115MB)} & \shortstack[c]{base \\ (366MB)} \\
                \midrule
                Move Playingcard   & 65 & \cellcolor{myblue}\textbf{87} & 72 \\
                Turn Switch        & 47 & 56 & \cellcolor{myblue}\textbf{64} \\
                Place Shoe         & 75 & 79 & \cellcolor{myblue}\textbf{85} \\
                Open Microwave     & 79 & \cellcolor{myblue}\textbf{98} & 92 \\
                Beat Block Hammer  & 74 & \cellcolor{myblue}\textbf{89} & 88 \\
                \midrule
                Average            & 68.0 & \cellcolor{myblue}\textbf{81.8} & 80.2 \\
                \bottomrule
            \end{tabular}
        }
    \end{minipage}
\end{table}
During the network design process, we systematically analyzed the impact of different encoder sizes and decoder depths on policy performance.

\textbf{Encoder Size}: We have already concluded that more advanced encoders with layer normalization, such as Uni3D\cite{zhou2023uni3d}, perform better than PointNet\cite{qi2017pointnet}. However, for the same type of encoder, using different ViT sizes as backbone results in significant differences in training efficiency and memory usage. Therefore, we investigated the impact of using different ViT sizes in the encoder on performance. Experimental results in Tab.~\ref{tab:encoder_size_comparison} and Tab.~\ref{tab:supptab1} show that selecting proper encoder sizes when dealing with different sizes of point clouds is crucial. We find that under the standard RoboTwin 2.0~\cite{chen2025robotwin} setting which parse the scene into 1024 points, ViT-tiny (53MB) is the optimal choice. While in real-world experiments (see section~\ref{sec:real_exp_detailed}) which uses 8192 points as input, a larger encoder such as ViT-small(115MB) is preferred. These results indicate that as the density of the point cloud increases, it is necessary to appropriately scale the ViT size within the visual encoder to ensure sufficient representational capacity for processing larger-scale point cloud data. Note that we pre-train each encoder using the corresponding number of points as input.

In Tab.~\ref{tab:supptab0}, under the 1024-point setting, we also scale the ViT sizes when dealing with different numbers of demonstrations for training. While we observe significant performance improvements across all encoder sizes as the number of policy training demonstrations increases, the ViT-tiny backbone consistently achieves the best results at 50, 100, and 200 demonstrations.

\textbf{Decoder Depth:}
We observed that the recent 3D policy ManiFlow\cite{yan2025maniflowgeneralrobotmanipulation} consistently uses a 12-layer DitX\cite{peebles2023scalable} decoder across all benchmarks and settings. However, networks with excessively deep layers can lead to severe overfitting. Therefore, we analyzed the impact of different decoder depths (in our method, it refers to the number of attention blocks detailed in Fig.~\ref{fig:pipe} b) on policy performance. As shown in Fig.~\ref{fig:decoder_depth}, when using the standard RoboTwin 2.0~\cite{chen2025robotwin} configuration (50 demonstrations for training), we found that decoders with 4 and 8 attention blocks showed significant performance advantages, and further increasing the decoder size beyond 8 blocks did not lead to performance improvements. Considering the trade-off between performance and efficiency, we adopted a decoder with 4 attention blocks in subsequent experiments. 

% \vspace{-2em}
% \section{Architecture}
% \label{sec:arch}
% % 先说最大的亮点：网络设计
% % architecture overview（左边dp3，右边是我们的）
% % 首先我们有一个更强的decoder结构，decoder结构是一个2way transformer adapt for action decode
% % 有一个分支，有一个cotraning task
% % besides，我们希望有一个有一个auxiliary task，compatible with different action space


% \vspace{}

\section{Experiment}
\label{sec: Exp}
% 上来1.先描述dataset的情况、多大训练集、相机、机器人
% 2.evaluation 测多少个episodes，是否沿用了之前的文章
\subsection{Simulation Experiments}
\begin{figure}[]
\vspace{-1.5em}
    \centering
    \includegraphics[width=\linewidth]{figures/sim5.pdf}
    \caption{\textbf{Simulation Settings:} The point clouds shown below are generated by cropping the table and the ground, followed by downsampling to 1024 points each.}
    \vspace{-0.5em}
    \label{fig:sim}
    \vspace{-1em}
\end{figure}
% \hspace{0.5ex}
\noindent \textbf{Simulation Environment.} We evaluate our method on two widely-used simulation benchmarks, RoboTwin 2.0~\cite{chen2025robotwin} and ManiSkill2~\cite{gu2023maniskill2}. RoboTwin 2.0~\cite{chen2025robotwin} is a scalable benchmark for evaluating robust bimanual manipulation across 50 tasks. It features an ``Easy'' setting with clean environments and a ``Hard'' setting that challenges policy robustness and generalization via strong domain randomization (e.g., random irrelevant objects, varying lighting and textures, larger pose variances). ManiSkill2~\cite{gu2023maniskill2} is a unified benchmark for manipulation skills, supported by a large-scale dataset of over 4 million demonstration frames.

\noindent \textbf{Experiment Setup.} For \textbf{RoboTwin 2.0}~\cite{chen2025robotwin}, we curated five distinct tasks for the ``Easy'' and ``Hard'' configurations, respectively. We follow the official setting to use 50 demos for each task using point clouds cropped to 1024 points. \\
For \textbf{ManiSkill2}~\cite{gu2023maniskill2}, we selected three representative tasks spanning both short and long horizon scenarios to evaluate the model's performance across varying temporal scales. For the ``PegInsertionSide'' task, we followed the method in Point Cloud Matters\cite{zhu2024point}, splitting the task into three subtasks: ``PegInsertion-Grasp'', ``PegInsertion-Align'', and ``PegInsertion-Insert'', which represent three states of this long-horizon task. For the ``PickCube'' task, we visualized the goal point as a green sphere within the point cloud for better representing of the goal position. For each ManiSkill2~\cite{gu2023maniskill2} task, we collected 1000 demonstrations with cropped point clouds (with the ground removed). In addition, we downsample point cloud to 1024 points using Farthest-Point-Sampling. Fig.~\ref{fig:sim} shows visualizations of the selected tasks and their corresponding point clouds. We report the success rate as the evaluation metric for these two benchmarks.

% \noindent \textbf{Baselines.} We evaluate our model against a comprehensive set of baselines on two benchmarks. For the Robotwin 2.0~\cite{chen2025robotwin} benchmark, we select representative methods across different input modalities.
% These include (1) 2D image-based policies: Diffusion Policy~\cite{chi2024diffusionpolicyvisuomotorpolicy} and RDT~\cite{liu2025rdt1bdiffusionfoundationmodel};
% (2) 2.5D RGBD-based policies: Spatial Forcing~\cite{li2025spatialforcingimplicitspatial}, DP+eVGGT~\cite{vuong2025improvingroboticmanipulationefficient}, and ACT+eVGGT~\cite{vuong2025improvingroboticmanipulationefficient};
% and (3) 3D pointcloud-based methods: 3D diffusion policy~\cite{ze20243d} and Maniflow~\cite{yan2025maniflowgeneralrobotmanipulation}. For the Maniskill2~\cite{gu2023maniskill2} benchmark, our comparison focuses on state-of-the-art 3D pointcloud-based methods.
% These are categorized into (1) methods trained from scratch: DP+SPuNet~\cite{zhu2024point}, DP+PointNet~\cite{zhu2024point}, 3D diffusion policy~\cite{ze20243d};
% and (2) methods utilizing 3D pretraining: ACT+PonderV2~\cite{zhu2024point}. \\


% dp3
% 再加一两个（sota）methods
% 一边找related work一边找一些baseline
% 最好选别人在这个benchmark上跑过的，可以直接抄过来，或者拿别人release的weights抄过来重跑
% fp3（github如果没有再升级过就不要去搞了）
% 搞不清楚的话可以让他们问问依阳
% 实在不行加dp
% --- 总的 figure/table 环境 ---
% \usepackage{tabularx} % 请确保在导言区添加此包

\begin{table*}[t]
\centering
\small
\captionsetup{justification=centering}
\caption{\textbf{Success Rate~(\%) on RoboTwin 2.0~\cite{chen2025robotwin} benchmark.}}
\label{tab: RoboTwin Benchmark}
\vspace{-1em}
% -------- 左侧子表：Easy (7行数据) --------
\begin{subtable}[t]{0.50\textwidth}
    \centering
    \caption{RoboTwin 2.0 ``Easy'' Setting}
    \vspace{-0.5em}
    \label{tab: Robotwin Easy}
    \resizebox{\linewidth}{!}{
        \renewcommand{\arraystretch}{1.0} % 左表保持标准行高
        \begin{tabular}{l ccccc c}
            \toprule
            \textbf{Method \textbackslash Task} &
            \shortstack{Beat \\ Hammer} &
            \shortstack{Move \\ Card} &
            \shortstack{Turn \\ Switch} &
            \shortstack{Open \\ Micro.} &
            \shortstack{Place \\ Shoe} &
            \textbf{Avg.} \\
            \midrule
            DP+eVGGT~\cite{vuong2025improvingroboticmanipulationefficient}              & 17 & 4  & 29 & 69 & 9  & 25.6 \\
            Spatial Forcing~\cite{li2025spatialforcingimplicitspatial}  & 47 & 60 & 37 & 85 & 22 & 50.2 \\
            DP3~\cite{ze20243d}        & 72 & 68 & 46 & 61 & 58 & 61.0 \\
            RDT~\cite{liu2025rdt1bdiffusionfoundationmodel}         & 77 & 43 & 35 & 37 & 35 & 45.4 \\
            DP~\cite{chi2024diffusionpolicyvisuomotorpolicy}          & 42 & 47 & 36 & 5  & 23 & 30.6 \\
            ManiFlow~\cite{yan2025maniflowgeneralrobotmanipulation}              & \cellcolor{myblue}\textbf{94} & 74 & 43 & 98 & 56 & 73.0 \\
            \textbf{Ours}         & 91 & \cellcolor{myblue}\textbf{85} & \cellcolor{myblue}\textbf{56} & \cellcolor{myblue}\textbf{100} & \cellcolor{myblue}\textbf{87} & \cellcolor{myblue}\textbf{83.8} \\
            \bottomrule
        \end{tabular}
    }
\end{subtable}
\hfill
% -------- 右侧子表：Hard (6行数据) --------
\begin{subtable}[t]{0.49\textwidth}
    \centering
    \caption{RoboTwin 2.0 ``Hard'' Setting}
    \vspace{-0.5em}
    \label{tab: RoboTwin Hard}
    \resizebox{\linewidth}{!}{
        % 关键：调大行间距以匹配左表高度 (7/6 ≈ 1.17，考虑到表头，1.2 左右最合适)
        \renewcommand{\arraystretch}{1.2} 
        \begin{tabular}{l ccccc c}
            \toprule
            \textbf{Method \textbackslash Task} &
            \shortstack{Beat \\ Hammer} &
            \shortstack{Lift \\ Pot} &
            \shortstack{Adjust \\ Bottle} &
            \shortstack{Burger \\ Fries} &
            \shortstack{Press \\ Stapler} &
            \textbf{Avg.} \\
            \midrule
            ACT+eVGGT~\cite{vuong2025improvingroboticmanipulationefficient}    & 7  & 2  & 20 & 10 & 14 & 10.6 \\
            DP+eVGGT~\cite{vuong2025improvingroboticmanipulationefficient}     & 18 & 14 & 52 & 26 & 20 & 26.0 \\
            DP3~\cite{ze20243d}          & 8  & 10 & 62 & 18 & 15 & 22.6 \\
            RDT~\cite{liu2025rdt1bdiffusionfoundationmodel}          & 9  & 12 & 50 & 20 & 20 & 22.2 \\
            ManiFlow~\cite{yan2025maniflowgeneralrobotmanipulation}         & 31 & 65 & 96 & 72 & \cellcolor{myblue}\textbf{44} & 61.6 \\
            \textbf{Ours}    & \cellcolor{myblue}\textbf{37} & \cellcolor{myblue}\textbf{67} & \cellcolor{myblue}\textbf{98} & \cellcolor{myblue}\textbf{95} & 27 & \cellcolor{myblue}\textbf{64.8} \\
            \bottomrule
        \end{tabular}
    }
\end{subtable}

% 在 subtable 下方添加注释说明
\vspace{-0.5em}
\end{table*}
% \end{table*}
\begin{table*}[t]
    \centering
    \small
    \captionsetup{justification=centering}
    \caption{\textbf{Success Rate~(\%) on ManiSkill2~\cite{gu2023maniskill2} benchmark.}}
    % \vspace{-0.5em}
    \label{tab: Maniskill2 results}
    
    \resizebox{\textwidth}{!}{
        % --- 关键修改：进一步调小行高倍数 ---
        \renewcommand{\arraystretch}{0.8} 
        \setlength{\tabcolsep}{8pt} 
        
        \begin{tabular}{l ccccc c} 
            \toprule
            \textbf{Method \textbackslash Task} &
            \shortstack{Pick \\ Cube} &
            \shortstack{Stack \\ Cube} &
            \shortstack{PegIns. \\ Grasp} &
            \shortstack{PegIns. \\ Align} &
            \shortstack{PegIns. \\ Insert} &
            \textbf{Average} \\
            % --- 也可以手动在首行后微调 ---
            \midrule
            DP+SPuNet~\cite{zhu2024point}       & 71 & 4  & 82 & 9  & 1 & 33.4 \\
            DP+PointNet~\cite{zhu2024point}      & 70 & 0  & 83 & 16 & 1 & 34.0 \\
            ACT+PonderV2~\cite{zhu2024point}     & 87 & \cellcolor{myblue}\textbf{35} & 65 & 23 & 2 & 42.4 \\
            DP3~\cite{ze20243d}                 & 57 & 14 & 41 & 3  & 0 & 23.0 \\
            \midrule
            \textbf{Ours}  & \cellcolor{myblue}\textbf{97} & 24 & \cellcolor{myblue}\textbf{97} & \cellcolor{myblue}\textbf{75} & \cellcolor{myblue}\textbf{21} & \cellcolor{myblue}\textbf{55.2} \\
            \bottomrule
        \end{tabular}
    }
    % --- 压缩表格与周围正文的间距 ---
    \vspace{-1em} 
\end{table*}




\noindent \textbf{Experimental Results.} As shown in Tab.~\ref{tab: RoboTwin Benchmark} and Tab.~\ref{tab: Maniskill2 results}, our method outperforms the 2D image-based, 2.5D RGBD-based, and 3D pointcloud-based baselines across two simulation benchmarks with following highlights:


\noindent \textbf{Generalization. } As illustrated in Fig. \ref{fig:sim}, the RoboTwin 2.0 ``Hard'' setting introduces significant challenges to models' generalization capabilities, featuring random distractors, massive objects, diverse textures, variable table heights, and lighting interference. Despite these complexities, our method demonstrates superior robustness. In Tab.~\ref{tab: RoboTwin Hard}, we achieve the highest success rates, demonstrating strong generalization ability in complex, unstructured environments.


\noindent \textbf{Fine-grained Manipulation} Task ``PegInsertion'' is challenging since it requires high spatial understanding of the environment to adjust the pose of the peg and accurate action generation to insert the peg into the hole. As shown in Tab. \ref{tab: Maniskill2 results}, while baselines suffer significant performance drops during the fine-grained Align and Insert stages (e.g., 23\% and 2\% for ACT+PonderV2), our method maintains success rates of 75\% and 21\%, respectively. Utilizing spatially aware dense features conditioning designs, our methods preserve the spatial resolution of 3D features, yielding better 3D understanding of fine-grained geometric details required for high-precision tasks.
% \vspace{-5pt}




\subsection{Ablation Study}\vspace{-1.5em}
\label{sec:ablation}
\begin{table*}[]
    \centering
    \small 
    \captionsetup{justification=centering} 
    \caption{\textbf{Ablation Study on Different Methods and Components~(\%).}}
    \label{tab:ablation_study_new}
    
    % 1. 使用 resizebox 强制限制宽度为 \textwidth
    % {!} 表示高度随宽度等比例缩放
    \resizebox{\textwidth}{!}{
        \renewcommand{\arraystretch}{1.0} % 进一步调小行高倍数
        \setlength{\tabcolsep}{6pt}      % 稍微恢复点列间距，因为 resizebox 会整体缩小文字
        
        \begin{tabular}{l c c c c c} 
            \toprule
            \textbf{Methods} & 
            \textbf{\shortstack[c]{Place \\ Shoe}} & 
            \textbf{\shortstack[c]{Move \\ Card.}} & 
            \textbf{\shortstack[c]{Beat \\ Hammer}} & 
            \textbf{\shortstack[c]{Turn \\ Switch}} & 
            \textbf{Average} \\
            \midrule
            
            Pretrained PointSAM (Global Feat. Cond.) & 60 & 76 & 86 & 28 & 62.50 \\
            
            Pretrained PointSAM (Dense Feat. Cond.) & \cellcolor{myblue}\textbf{92} & 70 & \cellcolor{myblue}\textbf{94} & 54 & 77.50 \\
            
            Train-from-Scratch PointSAM (Dense Feat. Cond.) & 60 & 58 & 90 & 55 & 65.75 \\
            
            \textbf{(Ours)} Pretrained-PointSAM (Dense) + Aux. EE & 87 & \cellcolor{myblue}\textbf{85} & 91 & \cellcolor{myblue}\textbf{56} & \cellcolor{myblue}\textbf{79.75} \\
            
            \bottomrule
        \end{tabular} 
    }
    \vspace{-3em}
\end{table*}
\begin{figure}[] % 使用 [t] 或 [htbp] 更好地控制位置
    \centering 
    % --- 第一个子图 (左侧) ---
    \begin{subfigure}[t]{0.48\textwidth} % 稍微调大比例，左右留一点间隙
        \centering
        \includegraphics[width=\textwidth]{figures/wandb1_new_cropped.pdf}
        \caption{Diffusion conditioned on global feature vs. attention-based spatially-aware dense feature.}
        \label{fig:diff_conditioning}
    \end{subfigure}
    \hfill % 使用 \hfill 撑开中间的空白，使两个子图分别靠左和靠右
    % --- 第二个子图 (右侧) ---
    \begin{subfigure}[t]{0.48\textwidth}
        \centering
        \includegraphics[width=\textwidth]{figures/wandb2_new_cropped.pdf}
        \caption{Pretrained encoder vs. train-from-scratch encoder.}
        \label{fig:pretrain}
    \end{subfigure}
% \vspace{-0.5em}
    % --- 总标题 ---
    \caption{Ablation study on different encoder design choices.} 
    \label{fig:combined_ablations}
     % 根据需要调整间距
\end{figure}
\vspace{-1.0em}
Here are some conclusions drawn from ablation studies (Tab.~\ref{tab:ablation_study_new} and Fig.~\ref{fig:combined_ablations}):
\begin{itemize}
    \item Replacing the global feature conditioning with our attention-based spatially aware dense feature conditioning yields a significant performance improvement (Tab.~\ref{tab:ablation_study_new}) and more stable training (Fig.~\ref{fig:diff_conditioning}).
    \item Encoder pretraining is effective, yielding better performance compared to training from scratch (Fig~\ref{fig:pretrain}).
    \item Incorporating auxiliary end-effector (EE) prediction yields a modest improvement in the average success rate.
\end{itemize}
\vspace{-10pt}
\subsection{Real-World Experiments}

\label{sec:real_exp_detailed}

% 【要点1：实验设计及原因（体现3D空间理解）】
% 这部分要说明为什么选择3D点云、为什么选择这些任务、为什么用多视角设置，体现对3D空间理解的重视
\begin{figure}[]
    \centering
    % 第一张图
    \begin{minipage}{0.35\textwidth}
        \centering
        \includegraphics[width=\linewidth]{figures/task_new3.pdf}
        \caption{Real-world setup.}
        \label{fig:setup}
        \vspace{-1em}
    \end{minipage}
    % 第二张图
    % \begin{minipage}{0.15\textwidth}
    %     \centering
    %     \includegraphics[width=\linewidth]{figures/disco_light_temp.pdf}
    %     \caption{Disco light.}
    %     \label{fig:disco}
    %     % \vspace{-2em}
    % \end{minipage}\hfill
    % 第三张图
    \begin{minipage}{0.6\textwidth}
        \centering
        \includegraphics[width=\linewidth]{figures/tasks222.pdf}
        \caption{Illustration of real-world settings and 3 real-world manipulation tasks.}
        \label{fig:execution_sequences}
        \vspace{-1em}
    \end{minipage}
\end{figure}
\textbf{Experimental Design for 3D Spatial Understanding.}
We further evaluate 3D Policies on a real-world robotic manipulation platform. We use an xArm6 robot arm with two Intel RealSense D435 RGB-D cameras (1 \textit{eye-to-hand} and 1 \textit{eye-in-hand}). We design 3 real-world tasks (see Fig.~\ref{fig:execution_sequences}) to validate the 3D understanding capabilities of policies. Specifically, they are a pose-constrained pick-and-place tasks—\textbf{\textit{(a)~Place Kettle on Stove}}—where successful execution entails grasping at the thin handle of the kettle and placing the kettle at designated locations with strictly defined spatial poses, an articulated-object manipulation task—\textbf{\textit{(b) Open Drawer}}, and a soft-body manipulation task—\textbf{\textit{(c) Fold Towel}}.
% Unlike 2D image-based methods that suffer from projective ambiguities and viewpoint dependencies, \textbf{3D point clouds provide a unified world-frame representation} that enables viewpoint invariance, geometric precision for sub-centimeter spatial reasoning, and natural multi-view fusion.
During the evaluation of each task, we introduce varied object poses (random position and orientation) to test the generalizability of all methods.

% \textbf{Experimental Setup.}
% We use a xArm6 robot arm with two Intel RealSense D435 RGB-D cameras (one eye-to-hand, one eye-in-hand). \textbf{3D point clouds provide a unified world-frame representation} enabling viewpoint invariance and geometric precision. We select six tasks: \textbf{\textit{(a) Place Cup}}, \textbf{\textit{(b) Fit Banana on Plate}}, \textbf{\textit{(c) Double Stack Block}}, \textbf{\textit{(d) Place Kettle on Stove}}, \textbf{\textit{(e) Open Drawer}}, \textbf{\textit{(f) Fold Towel}} (see Fig.~\ref{fig:execution_sequences}). Each task includes varied object poses (random position and orientation) to test the generalizability of policies.

% 【要点2：效果好】
% 展示实验结果表格和性能提升

% \textbf{Superior Performance Results.}
We collect 50 teleoperated demonstrations for training each task. Raw depth maps from sensors are enhanced using CDM~\cite{liu2025manipulation}; multi-view clouds are transformed to the robot base frame, cropped, and subsampled to 8,192 points via Farthest-Point-Sampling (FPS). We compare our method (using ViT-tiny, the same as simulation experiment settings) against baselines including DP3~\cite{ze20243d}, DP~\cite{chi2024diffusionpolicyvisuomotorpolicy}, ManiFlow~\cite{yan2025maniflowgeneralrobotmanipulation}, Pi0~\cite{black2024pi_0} using the exact same real-time fused point cloud input. Each data point is evaluated on 50 trials. See results in Tab.~\ref{tab:realworld_structured_transposed}.
\vspace{-1em}
% \begin{table*}[]
%     \centering
%     \small
%     \renewcommand{\arraystretch}{1.8}
%     \setlength{\tabcolsep}{8pt}
%     \caption{\textbf{Real-World Success Rates on Six Manipulation Tasks.}}
%     \label{tab:realworld}

%     \resizebox{0.95\textwidth}{!}{%
%     \begin{tabular}{l|c|ccc|c|c|c}
%         \toprule
%         \multirow{2}{*}{Model} &
%         \multirow{2}{*}{\raisebox{0.5em}{\shortstack{Pick \& Place\\(Long Horizon)}}} &
%         \multicolumn{3}{c|}{\shortstack{Pick \& Place}} &
%         \multirow{2}{*}{\raisebox{0.5em}{\shortstack{Articulation}}} &
%         \multirow{2}{*}{\raisebox{0.5em}{\shortstack{Soft\\Manipulation}}} &
%         \multirow{2}{*}{\raisebox{0.5em}{\textbf{Average (\%)}}} \\[0.5em]
%         \cmidrule(lr){2-2}
%         \cmidrule(lr){3-5}
%         \cmidrule(lr){6-6}
%         \cmidrule(lr){7-7}
%         & \shortstack{Double Stack\\Block} & \shortstack{Place\\Cup} & \shortstack{Fit Banana\\on Plate} & \shortstack{Place Kettle\\on Stove} & \shortstack{Open\\Drawer} & \shortstack{Fold\\Towel} & \\
%         \midrule
%         DP3           & 2 & 12 & 42 & 40 & 28 & 54 & 29.67 \\
%         DP            & 32 & 58 & 48 & 64 & 46 & 36 & 47.33 \\
%         \textbf{Ours} & \cellcolor{myblue}\textbf{12} & \cellcolor{myblue}\textbf{82} & \cellcolor{myblue}\textbf{64} & \cellcolor{myblue}\textbf{76} & \cellcolor{myblue}\textbf{64} & \cellcolor{myblue}\textbf{66} & \cellcolor{myblue}\textbf{60.67} \\
%         \bottomrule
%     \end{tabular}%
%     }
% \end{table*}
\begin{table}[] % 建议使用 [t] 或 [htbp]
    \centering
    \small
    \renewcommand{\arraystretch}{1.1}
    \setlength{\tabcolsep}{4pt}
    \caption{\textbf{Success Rates (\%) on Different Types of Real-World Tasks.}}
    \label{tab:realworld_structured_transposed}
    
    % 使用 \resizebox 强制将表格缩放到页面宽度 (\textwidth)
    % {!} 表示纵横比锁定，高度自动随宽度缩放
    \resizebox{\textwidth}{!}{
        \begin{tabular}{ll ccccc}
            \toprule
            \textbf{Task Type} & \textbf{Task}  & DP\cite{chi2024diffusionpolicyvisuomotorpolicy} & Pi0\cite{black2024pi_0} & DP3~\cite{ze20243d} & ManiFlow\cite{yan2025maniflowgeneralrobotmanipulation} & \textbf{Ours} \\
            \midrule
            Pick and Place & \textit{Place Kettle} & 64 & 46 & 40 & 36 & \cellcolor{myblue}\textbf{76} \\
            Articulated Object   & \textit{Open Drawer}  & 46 & 48 & 28 & 46 & \cellcolor{myblue}\textbf{64} \\
            Soft-body Object     & \textit{Fold Towel}   & 36 & 62 & 54 & 60 & \cellcolor{myblue}\textbf{66} \\
            \midrule
            &\textbf{Average} & \textbf{48.7} & \textbf{52.0} & \textbf{40.7} & \textbf{47.3} & \cellcolor{myblue}\textbf{68.7} \\
            \bottomrule
        \end{tabular}
    }
    \vspace{-1.1em}
\end{table}
% \vspace{-1em}
\begin{table}[htbp]
    \centering
    \small
    \renewcommand{\arraystretch}{1.2}
    \caption{\textbf{Point Cloud from Single Camera View vs. Two Camera Views.}}
    \label{tab:single_view}
    % 使用 resizebox 将表格宽度调整为页面宽度
    \resizebox{\textwidth}{!}{
        \setlength{\tabcolsep}{3.3pt}
        \begin{tabular}{lcccccccccc}
            \toprule
            \multirow{2}{*}{\textbf{Task}} & \multicolumn{2}{c}{DP~\cite{chi2024diffusionpolicyvisuomotorpolicy}} & \multicolumn{2}{c}{Pi0~\cite{black2024pi_0}} & \multicolumn{2}{c}{DP3~\cite{ze20243d}} & \multicolumn{2}{c}{ManiFlow~\cite{yan2025maniflowgeneralrobotmanipulation}} & \multicolumn{2}{c}{Ours} \\
            \cmidrule(lr){2-3} \cmidrule(lr){4-5} \cmidrule(lr){6-7} \cmidrule(lr){8-9} \cmidrule(lr){10-11}
            & Single & Two & Single & Two & Single & Two & Single & Two & Single & Two \\
            \midrule
            \textit{Place Kettle}    & 28 & 64 & - & 46 & 18 & 40 & 20 & 36 & \cellcolor{colorSingle}\textbf{40} & \cellcolor{colorTwo}\textbf{76} \\
            \textit{Open Drawer}     & 36 & 46 & - & 48 & 12 & 28 & 40 & 46 & \cellcolor{colorSingle}\textbf{52} & \cellcolor{colorTwo}\textbf{64} \\
            \textit{Fold Towel}      & 28 & 36 & - & 62 & 48 & 54 & 54 & 60 & \cellcolor{colorSingle}\textbf{62} & \cellcolor{colorTwo}\textbf{66} \\
            \midrule
            \textbf{Average} & \textbf{30.7} & \textbf{48.7} & \textbf{-} & \textbf{52.0} & \textbf{26.0} & \textbf{40.7} & \textbf{38.0} & \textbf{47.3} & \cellcolor{colorSingle}\textbf{51.3} & \cellcolor{colorTwo}\textbf{68.7} \\
            \bottomrule
        \end{tabular}
    }
    \vspace{-0.5em}
\end{table}

% As shown in Table~\ref{tab:realworld_structured}, our method achieves an average success rate of \textbf{60.67\% vs. 29.67\% (DP3)} and \textbf{60.67\% vs. 47.33\% (DP)}, demonstrating substantial improvements across all task types.


% 【要点3+4：原因分析（failure pattern和架构联系）+ 点云噪声处理】
% 分析DP3为什么失败，我们的架构为什么成功，把failure pattern和架构设计联系起来，同时说明点云噪声处理
\textbf{Number of Camera Views:} Previous 3D policies like DP3~\cite{ze20243d} and ManiFlow~\cite{yan2025maniflowgeneralrobotmanipulation} often rely on a single eye-to-hand camera setting, thereby \textbf{underutilizing the inherent capacity of point cloud representations to naturally fuse multi-view observations}. To further investigate, we evaluate the performance of all methods under two distinct settings: single-view (using only one eye-to-hand camera) and two-view configurations. Given that the pre-training of Pi0~\cite{black2024pi_0} relies on two views, only its two-view evaluation results are reported. Results are reported in Tab.~\ref{tab:single_view}, where all the methods show improvements when adding a second camera, and our method outperforms all baselines under both settings. In single-view experiments, we observe policies exhibit larger errors in tasks requiring fine-grained positional control (e.g., grasping a kettle handle or precisely aligning the kettle within a circular stove burner), and are prone to freezing mid-execution. Our spatially aware dense feature conditioning design not only yields better utilization of multi-view geometry, but also maintain robust performance with limited viewpoints.

\textbf{Robustness to external interference:} Point cloud color information is necessary for certain tasks
(e.g., placing a ``blue'' block on top of a ``red'' block). However, in real-world scenarios, point cloud colors are easily affected by changes of scene lighting at different times. To evaluate robustness to color variations, we introduce disco
light illumination (see Fig.~\ref{fig:setup}) that dynamically changes scene colors during evaluation, and test all methods under such environmental interference. In Tab.~\ref{tab:disco_light}, we observe performance degradation of all methods, however, our method maintains good performance and shows strong robustness under dynamic lighting conditions. 
\begin{table}[]
\vspace{-1em}
    \centering
    \small
    \renewcommand{\arraystretch}{1.2}
    \setlength{\tabcolsep}{3.9pt} 
    \caption{\textbf{Success Rate~(\%) under Normal and Disco Light.}}
    \label{tab:disco_light}
    \begin{tabular}{lcccccccccc}
        \toprule
        \multirow{2}{*}{\textbf{Task}} & \multicolumn{2}{c}{DP~\cite{chi2024diffusionpolicyvisuomotorpolicy}} & \multicolumn{2}{c}{Pi0~\cite{black2024pi_0}} & \multicolumn{2}{c}{DP3~\cite{ze20243d}} & \multicolumn{2}{c}{ManiFlow~\cite{yan2025maniflowgeneralrobotmanipulation}} & \multicolumn{2}{c}{Ours} \\
        \cmidrule(lr){2-3} \cmidrule(lr){4-5} \cmidrule(lr){6-7} \cmidrule(lr){8-9} \cmidrule(lr){10-11}
        & Nor. & Disc. & Nor. & Disc. & Nor. & Disc. & Nor. & Disc. & Nor. & Disc. \\
        \midrule
        \textit{Place Kettle} & 64 & 44 & 46 & 40 & 40 & 32 & 36 & 32 & \cellcolor{colorNormal}\textbf{76} & \cellcolor{colorDisco}\textbf{58} \\
        \textit{Open Drawer}  & 46 & 34 & 48 & 38 & 28 & 22 & 46 & 36 & \cellcolor{colorNormal}\textbf{64} & \cellcolor{colorDisco}\textbf{56} \\
        \textit{Fold Towel}   & 36 & 32 & 62 & 54 & 54 & 38 & 60 & 54 & \cellcolor{colorNormal}\textbf{66} & \cellcolor{colorDisco}\textbf{62} \\
        \midrule
        \textbf{Average} & \textbf{48.7} & \textbf{36.7} & \textbf{52.0} & \textbf{44.0} & \textbf{40.7} & \textbf{30.7} & \textbf{47.3} & \textbf{40.7} & \cellcolor{colorNormal}\textbf{68.7} & \cellcolor{colorDisco}\textbf{58.7} \\
        \bottomrule
    \end{tabular}
    % \vspace{-1em}
\end{table}
\vspace{-2em}
\section{Conclusion}
In conclusion, we address critical stability bottlenecks in 3D policy learning by proposing a robust transformer-diffusion architecture. Our approach significantly outperforms state-of-the-art baselines, establishing a scalable and effective foundation for generalizable 3D imitation learning.
\clearpage

% ---- Bibliography ----
%
% BibTeX users should specify bibliography style 'splncs04'.
% References will then be sorted and formatted in the correct style.
%
\bibliographystyle{splncs04}
\bibliography{main}
\end{document}
